\section[E. 右灯左行]{\myboxv{      \BF{\TT{E}}} 右灯左行}

\textit{运行时间限制}   1 second

\textit{内存使用限制} 1 Gigabyte

\textit{读入文件} \qquad 标准输入

\textit{输出文件} \qquad 标准输出


\begin{quote}
	\textit{的冲突，}
\end{quote}

\BF{本题和试题 \myboxiv{\BF{\TT{D}}} 左灯右行共用相同的背景，两道题之间有较大的关系，建议在尝试其中一题前首先阅读另一题的文本。}

\BF{这不是一道交互题。}

\subsection{题目描述}

小 T 面对 T 城的地铁图犯了难 ...

具体而言，这座城市在 E 河两岸，北岸有 $n$ 个地铁站点 $N_1$, $N_2$, $\ldots$, $N_n$, 南岸有 $n$ 个地铁站点 $S_1$, $S_2$, $\ldots$, $S_n$.

T 城总共有 $n - 1$ 条地铁线路，每条线路都是环线，其中第 $i$ 条环线连接的地铁站是 $N_i - N_{i + 1} - S_{i + 1} - S_i - N_i$, 但这些地铁环线都是单向运行的，可能是 $N_i \leftarrow N_{i + 1} \leftarrow S_{i + 1} \leftarrow S_i \leftarrow N_i$ 也可能是 $N_i \rightarrow N_{i + 1} \rightarrow S_{i + 1} \rightarrow S_i \rightarrow N_i$.


每条环线都有一个价目表 $u_i$, $d_i$, $l_i$, $r_i$, 分别表示这 4 个相邻站 ($N_i - N_{i + 1}$, $S_i - S_{i + 1}$, $N_i - S_i$ 和 $N_{i + 1} - S_{i + 1}$) 之间坐一站地铁移动的价格，注意坐一站地铁的方向由地铁的运行方向而定，可能是 $N_i \rightarrow N_{i + 1}$ 也可能是 $N_{i} \leftarrow N_{i + 1}$.

高兴的是，时空白洞吐出了价目表上这些环线的运行方向和价格，大 K 准备求助小 T 来得到一些通过地铁移动的价格。

大 K 会告诉小 T 这些环线的运行方向和价目表，随后小 T 需要回答 $q$ 个问题：

\begin{itemize}
  \item  地铁站 $x$(可以是 $N_i$ 也可以是 $S_i$, 下略) 到地铁站 $y$ 坐地铁，在给定的地铁价目表下，至少需要支付多少地铁费用？
  \item 若不同环线连接同一对相邻站，则这些边同时存在并且都可以使用。一次旅行的费用为所经过边权之和，站点间换乘不额外收费。
\end{itemize}

请帮助小 T 解决这个问题。

\subsection{输入格式}

第一行输入一个整数 $T$ ($ 1 \le T \le 10^5$), 表示测试数据的数量。

每组测试数据第一行输入两个整数 $n$ 和 $q$ ($2 \le n \le 10^5, 1 \le q \le 10^5$), 表示参数和询问数量，其中地铁站共 $2n$ 座，环线共 $n - 1$ 条。

随后输入 $n - 1$ 行，其中第 $i$ 行输入 $4$ 个整数 $u_i$, $d_i$, $l_i$, $r_i$($0 \le u_i, d_i, l_i, r_i \le 10^9$), 分别表示第 $i$ 条环线上 $N_i - N_{i + 1}$, $S_i - S_{i + 1}$, $N_i - S_i$ 和 $N_{i + 1} - S_{i + 1}$ 这四个相邻站之间的移动价格。

随后输入一行一个长度为 $n - 1$ 的字符串 $S$, 其中第 $i$ 个字符 $S_i \in \{\mathtt{I}, \mathtt{O}\}$, 若 $S_i = \mathtt{I}$ 则表示第 $i$ 条环线的方向为 $N_i \leftarrow N_{i + 1} \leftarrow S_{i + 1} \leftarrow S_i \leftarrow N_i$, 若 $S_i = \mathtt{O}$ 则表示第 $i$ 条环线的方向为 $N_i \rightarrow N_{i + 1} \rightarrow S_{i + 1} \rightarrow S_i \rightarrow N_i$.

随后输入 $q$ 行每行输入 $4$ 个元素 $S_x$, $id_x$, $S_y$, $id_y$, ($S_x, S_y \in \{\mathtt{N}, \mathtt{S}\}$, $id_x, id_y \in \{m \in \mathbb{N} : 1 \le m \le n\}$), 表示大 K 询问小 T 位于 $S_x$ 岸的第 $id_x$ 个地铁站 $(S_x)_{id_x}$ 到位于 $S_y$ 岸的第 $id_y$ 个地铁站 $(S_y)_{id_y}$ 之间至少需要支付多少地铁费用。

保证所有测试数据输入中的 $\sum n \le 3 \times 10^5$, $\sum q \le 3 \times 10^5$.

\subsection{输出格式}

每组测试数据输出 $q$ 行，第 $i$ 行输出一个整数 $d$ 表示小 T 收到的第 $i$ 个询问的答案，也即 $(S_x)_{id_x}$ 到 $(S_y)_{id_y}$ 至少需要支付多少地铁费用。


\subsection{样例 1}

\begin{minipage}{\linewidth}
	\begin{minipage}[t]{0.49\linewidth}
		\begin{lstlisting}[title=\texttt{standard input}]
3
3 5
10 3 1 2
10 10 1 1
OI
N 1 N 2
N 2 N 3
N 2 S 2
S 2 N 2
N 1 N 1
6 1
5 1 7 3
2 6 3 4
6 8 1 7
10 0 7 8
4 5 4 9
IOIOI
N 4 S 4
6 1
4 5 7 6
5 0 9 1
1 5 1 5
10 4 7 1
2 1 3 3
OIOIO
S 6 S 4
\end{lstlisting}
	\end{minipage}
	\begin{minipage}[t]{0.49\linewidth}
		\begin{lstlisting}[title=\texttt{standard output}]
10
12
1
14
0
14
17
\end{lstlisting}
	\end{minipage}
\end{minipage}


\subsection{样例 1 解释}

\begin{figure}[htbp]
   \centering
\begin{tikzpicture}[
  >=Stealth,
  station/.style={circle,draw,thick,minimum size=10mm,font=\large},
  bank/.style={thick},
  solid/.style={thick,->},
  rdashed/.style={thick,dashed,->},
  edgecost/.style={font=\Large, inner sep=0pt}
]

% --- coordinates / stations ---
\node[station] (N1) at (0, 2.2) {$N_1$};
\node[station] (N2) at (4, 2.2) {$N_2$};
\node[station] (N3) at (8, 2.2) {$N_3$};

\node[station] (S1) at (0,-2.2) {$S_1$};
\node[station] (S2) at (4,-2.2) {$S_2$};
\node[station] (S3) at (8,-2.2) {$S_3$};

% --- banks (text NOT on the line) ---
\draw[bank] (-1.5, 0.9) -- (9.5, 0.9);
\node[font=\Large] at (9.4, 1.25) {North bank};

\draw[bank] (-1.5,-0.9) -- (9.5,-0.9);
\node[font=\Large] at (9.4,-0.55) {South bank};

% --- Ring 1 = O : N1->N2->S2->S1->N1 ---
\draw[solid] (N1) -- (N2) node[midway,above=8pt,edgecost] {10};
\draw[solid] (N2) -- (S2) node[midway,left=10pt,edgecost] {2};


\draw[solid] (S2) -- (S1) node[midway,below=8pt,edgecost] {3};
\draw[solid] (S1) -- (N1) node[midway,left=10pt,edgecost] {1};

% --- Ring 2 = I : N3->N2, S3->N3, S2->S3, N2->S2 ---
\draw[rdashed] (N3) -- (N2) node[midway,above=8pt,edgecost] {10};
\draw[rdashed] (S3) -- (N3) node[midway,right=10pt,edgecost] {1};
\draw[rdashed] (S2) -- (S3) node[midway,below=8pt,edgecost] {10};

% --- the requested "double edge" on N2->S2: one solid + one dashed ---
% dashed copy slightly shifted to the right, so both are visible
\draw[rdashed] ([xshift=3pt]N2.south) -- ([xshift=3pt]S2.north)
  node[midway,right=14pt,edgecost] {1};

% --- ring labels ---
\node[font=\Large] at (2.0, 0) {Ring 1};
\node[font=\Large] at (6.0, 0) {Ring 2};

\end{tikzpicture}
\end{figure}

上面的图展示了第 1 组测试数据中地铁的运行方向和价目表。






